\documentclass[11pt]{article}
\usepackage[margin=0.82in]{geometry}
\usepackage{amsmath,amssymb,mathtools}
\usepackage{booktabs,tabularx,array}
\usepackage[dvipsnames]{xcolor}
\usepackage[colorlinks=true,linkcolor=NavyBlue,urlcolor=BrickRed,hypertexnames=false]{hyperref}
\usepackage{microtype}
\usepackage{fancyhdr}
\usepackage{enumitem}

\definecolor{ink}{HTML}{162235}
\definecolor{accent}{HTML}{245B78}
\definecolor{pale}{HTML}{EEF5F7}
\definecolor{passgreen}{HTML}{16794A}
\color{ink}
\pagestyle{fancy}
\fancyhf{}
\lhead{Restricted monomial groups}
\rhead{Proof and verification}
\cfoot{\thepage}
\setlength{\headheight}{14pt}
\newcommand{\Sfun}{\mathcal S}
\newcommand{\Cyc}{\operatorname{Cyc}}
\newcommand{\Sym}{\operatorname{Sym}}
\newcommand{\PASS}{\textcolor{passgreen}{\textbf{PASS}}}
\newcommand{\summarybox}[1]{\begin{center}\fcolorbox{accent}{pale}{\parbox{0.91\linewidth}{#1}}\end{center}}

\title{\textbf{Restricted Monomial Groups}\\[3pt]
\large Exact sigma--Molien formulas, proof, and matrix verification}
\author{Computational verification note}
\date{17 July 2026}

\begin{document}
\maketitle

\summarybox{\textbf{Outcome.}  The general dual-code formula and the specialized
$G(r,p,n)$ coefficient formula were compared with explicit matrix averages in
six representative cases.  All 180 coefficients through total degree six
agree.  A direct Reynolds projector for $G(2,2,3)$ on
$\Sym^2(V)\otimes V$ has rank and trace $1$, with idempotence error
$3.21\times10^{-16}$.}

\section{Statement of the problem}

Let $G$ be a finite group and $\rho:G\to GL(V)$ a complex representation.
For independent degree markers $\mathbf t=(t_1,t_2,\ldots)$ define
\begin{equation}
\Sfun_{G,V}(\mathbf t)
=\frac1{|G|}\sum_{g\in G}\prod_{a\geq1}
\det(I-t_a\rho(g))^{-1}.                                      \tag{1}
\end{equation}
The coefficient of $t_1^{\tau_1}\cdots t_\ell^{\tau_\ell}$ is claimed to be
\begin{equation}
\dim\left(\bigotimes_{a=1}^{\ell}\Sym^{\tau_a}V\right)^G.   \tag{2}
\end{equation}
We prove this master identity, specialize it to code-constrained monomial
groups, and independently test the resulting formulas on explicitly
constructed matrices.

\section{The master identity}

Suppose $g$ has eigenvalues $\lambda_1,\ldots,\lambda_n$.  The generating
function for complete homogeneous symmetric polynomials gives
\[
\det(I-tg)^{-1}=\prod_{i=1}^n(1-t\lambda_i)^{-1}
=\sum_{d\geq0}h_d(\lambda_1,\ldots,\lambda_n)t^d.
\]
But $h_d(\lambda_1,\ldots,\lambda_n)$ is the character of $\Sym^dV$ at $g$.
Consequently, the requested coefficient of (1) is
\[
\frac1{|G|}\sum_{g\in G}\prod_a\chi_{\Sym^{\tau_a}V}(g)
=\operatorname{tr}\left(\frac1{|G|}\sum_{g\in G}
g\ \middle|\ \bigotimes_a\Sym^{\tau_a}V\right).
\]
The averaged operator is the Reynolds projector onto the invariant subspace.
Its trace is its rank, proving (2).  This proof also explains the three
computational routes used below: determinant coefficients, character averages,
and projector ranks must agree.

\section{Code-constrained monomial groups}

Fix $\zeta=e^{2\pi i/r}$, an $H$-stable additive code
$A\leq(\mathbb Z/r\mathbb Z)^n$, and
\[
G(A,H)=D(A)\rtimes H,\qquad
D(A)=\{\operatorname{diag}(\zeta^{x_1},\ldots,\zeta^{x_n}):x\in A\}.
\]
We use the convention $g=D(x)P_\pi$.  If $C$ is a cycle of $\pi$, the
restriction of $g$ to its coordinate subspace is a weighted cyclic shift.
Only the full traversal contributes to its characteristic polynomial, hence
\begin{equation}
\det(I-tg)=\prod_{C\in\Cyc(\pi)}
\left(1-\zeta^{x_C}t^{|C|}\right),\qquad
x_C=\sum_{i\in C}x_i.                                         \tag{3}
\end{equation}
Substitution in (1) proves the finite exact formula
\begin{equation}
\Sfun_{A,H}(\mathbf t)=\frac1{|H||A|}
\sum_{\pi\in H}\sum_{x\in A}
\prod_{a,C}\left(1-\zeta^{x_C}t_a^{|C|}\right)^{-1}.         \tag{4}
\end{equation}

\subsection{The dual-code filter}

Expand each factor geometrically and let $m_{C,a}\geq0$ be its selected
multiplicity.  Put $M_C=\sum_am_{C,a}$ and set $e_i=M_C$ for $i\in C$.  The
phase of a term is $\zeta^{x\cdot e}$.  For
\[
A^\perp=\{e\in(\mathbb Z/r)^n:x\cdot e=0\pmod r\text{ for all }x\in A\},
\]
finite character orthogonality gives
\[
\frac1{|A|}\sum_{x\in A}\zeta^{x\cdot e}=\mathbf1_{e\in A^\perp}.
\]
Therefore the coefficient indexed by $\tau=(\tau_1,\ldots,\tau_\ell)$ is
\begin{equation}
\frac1{|H|}\sum_{\pi\in H}
\#\left\{(m_{C,a}):
\begin{array}{l}
\sum_C|C|m_{C,a}=\tau_a\quad\text{for every }a,\\
e(\pi,m)\in A^\perp
\end{array}\right\}.                                        \tag{5}
\end{equation}
Equation (5) is the implemented general formula.  Notice that it contains no
matrix eigenvalue computation and no average over diagonal matrices.

For a residue $s\in\mathbb Z/r$, the equivalent root-of-unity filter is
\[
\Phi_{d,s}(\mathbf t)=\frac1r\sum_{\ell=0}^{r-1}
\zeta^{-\ell s}\prod_a(1-\zeta^\ell t_a^d)^{-1}.
\]
It collects those monomials for which $\sum_am_a\equiv s\pmod r$.

\section{Specialization to $G(r,p,n)$}

Assume $p\mid r$ and
\[
A_{r,p,n}=\{x\in(\mathbb Z/r)^n:\sum_i x_i=0\pmod p\}.
\]
The vectors $e_i-e_j\in A_{r,p,n}$ force a dual vector to have equal
coordinates.  The vectors $p e_i\in A_{r,p,n}$ then force the common coordinate
to be a multiple of $r/p$.  Thus
\begin{equation}
A_{r,p,n}^\perp
=\{q(r/p)(1,\ldots,1):0\leq q<p\}.                            \tag{6}
\end{equation}
The cycle-index theorem now yields
\begin{equation}
\Sfun_{G(r,p,n)}(\mathbf t)
=\sum_{q=0}^{p-1}[u^n]\exp\left(
\sum_{d\geq1}\frac{u^d}{d}\Phi_{d,qr/p}(\mathbf t)\right).  \tag{7}
\end{equation}
An equivalent coefficient algorithm regards each coordinate as carrying a
label $(m_1,\ldots,m_\ell)$ whose coordinate sum has residue $qr/p$.  An
$S_n$-orbit is a multiset of $n$ labels.  Bounded dynamic programming over
label multiplicities extracts exactly the coefficient of (7); this is the
second independent formula route in the software.

With only $t_1=t$ nonzero, coefficient extraction reduces to
\begin{equation}
\Sfun_{G(r,p,n)}(t,0,\ldots)=
\frac1{(1-t^{nr/p})\prod_{j=1}^{n-1}(1-t^{rj})}.              \tag{8}
\end{equation}
For $W(D_n)=G(2,2,n)$, the second residue sector requires every coordinate to
carry odd total multiplicity.  It cannot occur below degree $n$.  In
particular, it does not create a degree-one fixed vector; this is the key
zero-coordinate correction in the proposed discussion.

\section{Independent verification design}

\begin{enumerate}[leftmargin=*,itemsep=3pt]
\item Construct every matrix $D(x)P_\pi$.  Check the predicted order
$r^n n!/p$ in the $G(r,p,n)$ cases.
\item For every partition $\tau$ of every total degree from $0$ through $6$,
compute $h_{\tau_a}$ recursively from the eigenvalues of each matrix and
average $\prod_a h_{\tau_a}$ over the group.
\item Independently evaluate (5) for a general code case or the multiset
coefficient of (7) for $G(r,p,n)$.
\item For a selected target, explicitly embed normalized occupancy bases of
symmetric powers in tensor powers, form the full Reynolds matrix, and compare
its numerical rank and trace with the coefficient.
\end{enumerate}

These routes share only the mathematical input $(r,p,n)$ or $(A,H)$.  The
formula route does not call the matrix-average code.

\section{Representative cases and results}

There are 30 partitions through total degree six.  The following suite includes
an $H=C_3$ code group not covered by the full symmetric-group specialization,
the unrestricted wreath products, the $D_n$ parity sector, and two different
nontrivial divisors $p$ of $r$.

\begin{center}
\begin{tabularx}{\linewidth}{@{}Xrrrrc@{}}
\toprule
Case & $|G|$ & $\dim V$ & degree & checks & result\\
\midrule
even-sign code $\rtimes C_3$ & 12 & 3 & $0$--$6$ & 30 & \PASS\\
$G(2,1,2)$ & 8 & 2 & $0$--$6$ & 30 & \PASS\\
$G(2,2,3)=W(D_3)$ & 24 & 3 & $0$--$6$ & 30 & \PASS\\
$G(3,1,2)$ & 18 & 2 & $0$--$6$ & 30 & \PASS\\
$G(4,2,2)$ & 16 & 2 & $0$--$6$ & 30 & \PASS\\
$G(4,4,2)$ & 8 & 2 & $0$--$6$ & 30 & \PASS\\
\bottomrule
\end{tabularx}
\end{center}

Selected matched dimensions illustrate that the test is multivariate, not only
the ordinary Molien specialization.
\begin{center}
\begin{tabular}{@{}lrrrr@{}}
\toprule
Case & $(2)$ & $(3)$ & $(2,1)$ & $(3,2)$\\
\midrule
even-sign code $\rtimes C_3$ & 1 & 1 & 1 & 4\\
$G(2,1,2)$ & 1 & 0 & 0 & 0\\
$G(2,2,3)$ & 1 & 1 & 1 & 3\\
$G(3,1,2)$ & 0 & 1 & 1 & 0\\
$G(4,2,2)$ & 0 & 0 & 0 & 0\\
$G(4,4,2)$ & 1 & 0 & 0 & 0\\
\bottomrule
\end{tabular}
\end{center}

For $G(2,2,3)$ and $\tau=(2,1)$, the target dimension is
$\binom{4}{2}\binom{3}{1}=18$.  Direct averaging of the eighteen-dimensional
target matrices produced
\[
\operatorname{rank}P=1,\qquad \operatorname{tr}P=0.9999999999999998,
\qquad\|P^2-P\|_F=3.21\times10^{-16}.
\]
This agrees with both coefficient routes.

\section{Reproducibility and interpretation}

Run the command-line suite from the repository root with
{\scriptsize
\begin{verbatim}
.venv/bin/python -m for_this_guy.matrix_group_formula_verification.restricted_monomial.verification
\end{verbatim}
}
or open \texttt{marimo\_notebooks/restricted\_monomial\_verification.py} with marimo.  Matrix
averages are accepted only when their imaginary part and distance from the
nearest integer are below $3\times10^{-7}$.  The observed residuals are much
smaller; the displayed projector error is at machine precision.

Finite tests do not replace the preceding proof.  Their role is to detect
indexing errors, wrong permutation conventions, missing residue sectors, and
incorrect treatment of zero-labeled coordinates.  Within the stated small
dimensions, no such discrepancy remains.

\end{document}
